% GENERATED FILE. Edit canonical lessons, then run npm run generate:books.
% canonical-source-hash: fnv1a64:ef0cc2c5da0b32a7
% canonical-lessons: MR-C461-apman, MR-C461-tomna, MR-C461-tiraskar, MR-C461-dukhavne, MR-C461-hinavne

\chapter{Hurt Feelings}
\label{ch:mr-hurt-feelings}
\hlchaptermodality{461}

\begin{chapteropening}
\textbf{By the end of this chapter:} \emph{I can talk about an insult, a taunt and hurt feelings.}

Complete the last lesson of chapter 461: I can talk about an insult, a taunt and hurt feelings.
\end{chapteropening}

\section[\texorpdfstring{\mr{अपमान} (apmān)}{अपमान (apmān)}]{\mr{अपमान} (apmān) — an insult; humiliation}

\label{lesson:MR-C461-apman}



\begin{quote}
Before the new one: say the Marathi for closeness, then the Marathi for beloved.
\end{quote}

\subsection*{You'll want to know: \mr{अपमान}}
\textbf{\mr{अपमान}} — \emph{apmān} — \textquotedblleft{}an insult; humiliation\textquotedblright{}.

\begin{quote}
\mr{सगळ्यांसमोर} \mr{झालेला} \mr{अपमान} \mr{तो} \mr{विसरला} \mr{नाही}.
\end{quote}

\begin{quote}
\emph{sagḷyāṁsamor jhālelā apmān to visarlā nāhī.}
\end{quote}

\begin{quote}
He did not forget being humiliated in front of everyone.
\end{quote}

Masculine.

\mr{अपमान} \mr{करणे}, to insult; \mr{अपमान} \mr{होणे}, to be humiliated.

\begin{cousinweb}
\textbf{\mr{अपमान}} (\emph{apmān}) is a tatsama, taken from Sanskrit \emph{apamāna}, \textquotedblleft{}disrespect, disgrace\textquotedblright{} (CDIAL 433).
\end{cousinweb}

\begin{grammarlens}[title={What you've built}]
One more word for describing what you lived and felt.
\end{grammarlens}

\subsection*{Your turn}
\begin{itemize}
\raggedright
  \item \emph{Say it:} \emph{apmān}
  \item \emph{Say it:} \emph{apmān}, once more
  \item \emph{Say it:} \emph{sagḷyāṁsamor jhālelā apmān to visarlā nāhī.} — He did not forget being humiliated in front of everyone.
  \item \emph{Recall:} say the Marathi for closeness, then the Marathi for beloved, then say \emph{apmān} again
\end{itemize}

\begin{tcolorbox}[breakable,colback=teal!4,colframe=teal!35!black,title={Before you move on}]
What does \mr{अपमान} mean? (\textquotedblleft{}An insult; humiliation\textquotedblright{}.) Say it once more.
\end{tcolorbox}
\section[\texorpdfstring{\mr{टोमणा} (ṭomṇā)}{टोमणा (ṭomṇā)}]{\mr{टोमणा} (ṭomṇā) — a taunt; a dig}

\label{lesson:MR-C461-tomna}



\begin{quote}
Before the new one: say the Marathi for beloved, then the Marathi for an insult.
\end{quote}

\subsection*{You'll want to know: \mr{टोमणा}}
\textbf{\mr{टोमणा}} — \emph{ṭomṇā} — \textquotedblleft{}a taunt; a dig\textquotedblright{}.

\begin{quote}
\mr{शेजारच्या} \mr{बाईंनी} \mr{आईला} \mr{टोमणा} \mr{मारला}.
\end{quote}

\begin{quote}
\emph{śejārcyā bāīnnī āīlā ṭomṇā mārlā.}
\end{quote}

\begin{quote}
The lady next door made a dig at Mother.
\end{quote}

Masculine.

\mr{टोमणा} \mr{मारणे}, to make a dig: the target takes -\mr{ला}.

\begin{cousinweb}
\textbf{\mr{टोमणा}} (\emph{ṭomṇā}) is a Marathi noun whose origin is not settled; CDIAL lists a Marathi \emph{ṭumṇī}, \textquotedblleft{}goading\textquotedblright{}, under \emph{trumpati} (CDIAL 6070), but does not list this word.
\end{cousinweb}

\begin{grammarlens}[title={What you've built}]
One more word for describing what you lived and felt.
\end{grammarlens}

\subsection*{Your turn}
\begin{itemize}
\raggedright
  \item \emph{Say it:} \emph{ṭomṇā}
  \item \emph{Say it:} \emph{ṭomṇā}, once more
  \item \emph{Say it:} \emph{śejārcyā bāīnnī āīlā ṭomṇā mārlā.} — The lady next door made a dig at Mother.
  \item \emph{Recall:} say the Marathi for beloved, then the Marathi for an insult, then say \emph{ṭomṇā} again
\end{itemize}

\begin{tcolorbox}[breakable,colback=teal!4,colframe=teal!35!black,title={Before you move on}]
What does \mr{टोमणा} mean? (\textquotedblleft{}A taunt; a dig\textquotedblright{}.) Say it once more.
\end{tcolorbox}
\section[\texorpdfstring{\mr{तिरस्कार} (tiraskār)}{तिरस्कार (tiraskār)}]{\mr{तिरस्कार} (tiraskār) — contempt}

\label{lesson:MR-C461-tiraskar}



\begin{quote}
Before the new one: say the Marathi for an insult, then the Marathi for a taunt.
\end{quote}

\subsection*{You'll want to know: \mr{तिरस्कार}}
\textbf{\mr{तिरस्कार}} — \emph{tiraskār} — \textquotedblleft{}contempt\textquotedblright{}.

\begin{quote}
\mr{त्याला} \mr{खोटेपणाचा} \mr{तिरस्कार} \mr{वाटतो}.
\end{quote}

\begin{quote}
\emph{tyālā khoṭepaṇācā tiraskār vāṭto.}
\end{quote}

\begin{quote}
He feels contempt for dishonesty.
\end{quote}

Masculine.

\mr{तिरस्कार} \mr{वाटणे}, to feel contempt, with the dative; \mr{तिरस्कार} \mr{करणे}, to scorn.

\begin{cousinweb}
\textbf{\mr{तिरस्कार}} (\emph{tiraskār}) is a tatsama, taken from Sanskrit \emph{tiraskāra}, \textquotedblleft{}placing aside, concealment\textquotedblright{} (Monier-Williams); in Marathi it is contempt.
\end{cousinweb}

\begin{grammarlens}[title={What you've built}]
One more word for describing what you lived and felt.
\end{grammarlens}

\subsection*{Your turn}
\begin{itemize}
\raggedright
  \item \emph{Say it:} \emph{tiraskār}
  \item \emph{Say it:} \emph{tiraskār}, once more
  \item \emph{Say it:} \emph{tyālā khoṭepaṇācā tiraskār vāṭto.} — He feels contempt for dishonesty.
  \item \emph{Recall:} say the Marathi for an insult, then the Marathi for a taunt, then say \emph{tiraskār} again
\end{itemize}

\begin{tcolorbox}[breakable,colback=teal!4,colframe=teal!35!black,title={Before you move on}]
What does \mr{तिरस्कार} mean? (\textquotedblleft{}Contempt\textquotedblright{}.) Say it once more.
\end{tcolorbox}
\section[\texorpdfstring{\mr{दुखावणे} (dukhāvṇe)}{दुखावणे (dukhāvṇe)}]{\mr{दुखावणे} (dukhāvṇe) — to hurt someone's feelings}

\label{lesson:MR-C461-dukhavne}



\begin{quote}
Before the new one: say the Marathi for a taunt, then the Marathi for contempt.
\end{quote}

\subsection*{You'll want to know: \mr{दुखावणे}}
\textbf{\mr{दुखावणे}} — \emph{dukhāvṇe} — \textquotedblleft{}to hurt someone's feelings\textquotedblright{}.

\begin{quote}
\mr{माझ्या} \mr{बोलण्याने} \mr{मी} \mr{तिला} \mr{दुखावलं}.
\end{quote}

\begin{quote}
\emph{mājhyā bolṇyāne mī tilā dukhāvlaṁ.}
\end{quote}

\begin{quote}
I hurt her feelings with what I said.
\end{quote}

Past, with -\mr{ने} (he …): \mr{त्याने} \mr{तिला} \mr{दुखावलं}, \emph{tyāne tilā dukhāvlaṁ}.

The one hurt takes -\mr{ला}; \mr{मन} \mr{दुखावलं}, feelings were hurt.

\begin{cousinweb}
\textbf{\mr{दुखावणे}} (\emph{dukhāvṇe}) is a tadbhava: CDIAL lists Marathi \emph{dukhāviṇẽ} under Sanskrit \emph{duḥkhati}, \textquotedblleft{}pains\textquotedblright{} (CDIAL 6376).
\end{cousinweb}

\begin{grammarlens}[title={What you've built}]
One more word for describing what you lived and felt.
\end{grammarlens}

\subsection*{Your turn}
\begin{itemize}
\raggedright
  \item \emph{Say it:} \emph{dukhāvṇe}
  \item \emph{Say it:} \emph{dukhāvṇe}, once more
  \item \emph{Say it:} \emph{mājhyā bolṇyāne mī tilā dukhāvlaṁ.} — I hurt her feelings with what I said.
  \item \emph{Recall:} say the Marathi for a taunt, then the Marathi for contempt, then say \emph{dukhāvṇe} again
\end{itemize}

\begin{tcolorbox}[breakable,colback=teal!4,colframe=teal!35!black,title={Before you move on}]
What does \mr{दुखावणे} mean? (\textquotedblleft{}To hurt someone's feelings\textquotedblright{}.) Say it once more.
\end{tcolorbox}
\section[\texorpdfstring{\mr{हिणवणे} (hiṇavṇe)}{हिणवणे (hiṇavṇe)}]{\mr{हिणवणे} (hiṇavṇe) — to belittle}

\label{lesson:MR-C461-hinavne}



\begin{quote}
Before the new one: say the Marathi for contempt, then the Marathi for to hurt someone's feelings.
\end{quote}

\subsection*{You'll want to know: \mr{हिणवणे}}
\textbf{\mr{हिणवणे}} — \emph{hiṇavṇe} — \textquotedblleft{}to belittle\textquotedblright{}.

\begin{quote}
\mr{दारिद्र्यावरून} \mr{कोणालाही} \mr{हिणवू} \mr{नये}.
\end{quote}

\begin{quote}
\emph{dāridryāvarūn koṇālāhī hiṇvū naye.}
\end{quote}

\begin{quote}
No one should be belittled for being poor.
\end{quote}

Past, with -\mr{ने} (he …): \mr{त्याने} \mr{तिला} \mr{हिणवलं}, \emph{tyāne tilā hiṇavlaṁ}.

The one belittled takes -\mr{ला}; the reason takes -\mr{वरून}.

\begin{cousinweb}
\textbf{\mr{हिणवणे}} (\emph{hiṇavṇe}) is a tadbhava: CDIAL lists Marathi \emph{hiṇaviṇẽ}, \textquotedblleft{}to lower in estimation\textquotedblright{}, under \emph{*hīnayati}, \textquotedblleft{}depreciates\textquotedblright{}, built on Sanskrit \emph{hīna}, \textquotedblleft{}low\textquotedblright{} (CDIAL 14126a).
\end{cousinweb}

\begin{grammarlens}[title={What you've built}]
One more word for describing what you lived and felt.
\end{grammarlens}

\subsection*{Your turn}
\begin{itemize}
\raggedright
  \item \emph{Say it:} \emph{hiṇavṇe}
  \item \emph{Say it:} \emph{hiṇavṇe}, once more
  \item \emph{Say it:} \emph{dāridryāvarūn koṇālāhī hiṇvū naye.} — No one should be belittled for being poor.
  \item \emph{Recall:} say the Marathi for contempt, then the Marathi for to hurt someone's feelings, then say \emph{hiṇavṇe} again
\end{itemize}

\begin{tcolorbox}[breakable,colback=teal!4,colframe=teal!35!black,title={Before you move on}]
What does \mr{हिणवणे} mean? (\textquotedblleft{}To belittle\textquotedblright{}.) Say it once more.
\end{tcolorbox}
